\documentclass[11pt,a4paper]{article}

\usepackage[utf8]{inputenc}
\usepackage[T1]{fontenc}
\usepackage{amsmath,amssymb,amsthm}
\usepackage{algorithm}
\usepackage{algpseudocode}
\usepackage{booktabs}
\usepackage{graphicx}
\usepackage{hyperref}
\usepackage{listings}
\usepackage{xcolor}
\usepackage{geometry}
\usepackage{tikz}
\usetikzlibrary{matrix,positioning,arrows.meta}
\geometry{margin=1in}

\newtheorem{theorem}{Theorem}
\newtheorem{lemma}[theorem]{Lemma}
\newtheorem{definition}{Definition}
\newtheorem{proposition}{Proposition}
\newtheorem{corollary}{Corollary}

\lstset{
    basicstyle=\ttfamily\small,
    breaklines=true,
    frame=single,
    numbers=left,
    numberstyle=\tiny,
    keywordstyle=\color{blue},
    commentstyle=\color{green!60!black},
    stringstyle=\color{red}
}

\title{Poseidon2 ZK-Friendly Hash Precompile for EVM:\\Enabling Efficient Zero-Knowledge Proofs}

\author{Zach Kelling\thanks{zach@lux.network} \\ \textit{Hanzo Industries \quad Lux Industries \quad Zoo Labs Foundation} \\ \texttt{research@lux.network}}

\date{June 2023}

\begin{document}

\maketitle

\begin{abstract}
We present an EVM precompiled contract for Poseidon2, a zero-knowledge-friendly cryptographic hash function designed for optimal performance in arithmetic circuits. Poseidon2 achieves significant improvements over its predecessor through a restructured round function that reduces constraint counts by 30-40\% while maintaining equivalent security margins. This precompile enables efficient on-chain verification of zero-knowledge proofs that utilize Poseidon2 as their hash primitive, supporting both Merkle tree constructions and general hashing operations. We provide a comprehensive gas cost model, analyze the algebraic security properties, and demonstrate integration patterns for zkSNARK and zkSTARK verification. Our implementation achieves Poseidon2 hashing at approximately 45,000 gas for two-element inputs over the BN254 scalar field, making zero-knowledge applications economically viable on EVM chains.
\end{abstract}

\section{Introduction}

Zero-knowledge proof systems have emerged as a transformative technology for blockchain scalability and privacy. These systems require cryptographic hash functions for commitment schemes, Merkle trees, and Fiat-Shamir transformations. Traditional hash functions like SHA-256 and Keccak are computationally expensive when expressed as arithmetic circuits, requiring thousands of constraints per hash evaluation.

Poseidon, introduced by Grassi et al.~\cite{poseidon}, addressed this by designing a hash function specifically optimized for prime field arithmetic. Poseidon2~\cite{poseidon2} further improves upon this design with a restructured round function that achieves:

\begin{itemize}
    \item 30-40\% fewer constraints in R1CS/Plonk circuits
    \item Improved native performance for verification
    \item Equivalent security margins
    \item Better parallelization potential
\end{itemize}

This paper presents LIP-003, specifying an EVM precompile for Poseidon2 with:

\begin{enumerate}
    \item Support for multiple prime fields (BN254, BLS12-381 scalar fields)
    \item Configurable width and rate parameters
    \item Sponge construction for arbitrary-length inputs
    \item Merkle tree hashing optimizations
    \item Comprehensive security analysis
\end{enumerate}

\section{Poseidon2 Design}

\subsection{Algebraic Hash Functions}

\begin{definition}[Algebraic Hash Function]
A hash function $H: \mathbb{F}_p^* \to \mathbb{F}_p^n$ is algebraic if it can be expressed as a sequence of additions, multiplications, and constant operations over a finite field $\mathbb{F}_p$.
\end{definition}

Algebraic hash functions achieve efficiency in zero-knowledge circuits because:
\begin{enumerate}
    \item Field operations are native to the circuit arithmetic
    \item No bit decomposition required (unlike SHA/Keccak)
    \item Low multiplicative depth enables efficient proof generation
\end{enumerate}

\subsection{Poseidon2 Round Function}

Poseidon2 operates on a state of $t$ field elements through $R = R_F + R_P$ rounds:
\begin{itemize}
    \item $R_F/2$ full rounds at beginning and end
    \item $R_P$ partial rounds in the middle
\end{itemize}

\begin{algorithm}
\caption{Poseidon2 Permutation}
\label{alg:poseidon2}
\begin{algorithmic}[1]
\Require State $\mathbf{s} \in \mathbb{F}_p^t$, round constants $\mathbf{c}$, matrices $M_E$, $M_I$
\Ensure Permuted state $\mathbf{s}'$
\State $\mathbf{s} \gets M_E \cdot \mathbf{s}$ \Comment{Initial external matrix}
\For{$r = 0$ to $R_F/2 - 1$} \Comment{First full rounds}
    \State $\mathbf{s} \gets \mathbf{s} + \mathbf{c}_r$ \Comment{Add round constants}
    \State $\mathbf{s} \gets \mathbf{s}^{\alpha}$ \Comment{Full S-box layer}
    \State $\mathbf{s} \gets M_E \cdot \mathbf{s}$ \Comment{External matrix}
\EndFor
\For{$r = R_F/2$ to $R_F/2 + R_P - 1$} \Comment{Partial rounds}
    \State $s_0 \gets s_0 + c_r$ \Comment{Single constant}
    \State $s_0 \gets s_0^{\alpha}$ \Comment{Single S-box}
    \State $\mathbf{s} \gets M_I \cdot \mathbf{s}$ \Comment{Internal matrix}
\EndFor
\For{$r = R_F/2 + R_P$ to $R_F + R_P - 1$} \Comment{Final full rounds}
    \State $\mathbf{s} \gets \mathbf{s} + \mathbf{c}_r$
    \State $\mathbf{s} \gets \mathbf{s}^{\alpha}$
    \State $\mathbf{s} \gets M_E \cdot \mathbf{s}$
\EndFor
\Return $\mathbf{s}$
\end{algorithmic}
\end{algorithm}

\subsection{Key Improvements over Poseidon}

\begin{enumerate}
    \item \textbf{Split matrix design}: External matrix $M_E$ for full rounds, simpler internal matrix $M_I$ for partial rounds
    \item \textbf{Efficient internal matrix}: $M_I = I + \mathbf{1}\mathbf{\mu}^T$ where $\mathbf{\mu}$ has special structure
    \item \textbf{Reduced constraints}: Internal matrix multiplication requires only $t$ additions
    \item \textbf{Better diffusion}: External matrix provides stronger mixing in full rounds
\end{enumerate}

\subsection{Matrix Structures}

\textbf{External Matrix $M_E$} (for full rounds):
\[
M_E = \text{circ}(M_4) \text{ where } M_4 = \begin{pmatrix} 2 & 3 & 1 & 1 \\ 1 & 2 & 3 & 1 \\ 1 & 1 & 2 & 3 \\ 3 & 1 & 1 & 2 \end{pmatrix}
\]

\textbf{Internal Matrix $M_I$} (for partial rounds):
\[
M_I = \begin{pmatrix} \mu_0 + 1 & \mu_1 & \cdots & \mu_{t-1} \\ \mu_0 & \mu_1 + 1 & \cdots & \mu_{t-1} \\ \vdots & \vdots & \ddots & \vdots \\ \mu_0 & \mu_1 & \cdots & \mu_{t-1} + 1 \end{pmatrix}
\]

This structure allows $M_I \cdot \mathbf{s}$ to be computed as:
\[
(M_I \cdot \mathbf{s})_i = s_i + \sum_{j=0}^{t-1} \mu_j \cdot s_j
\]

\section{Precompile Specification}

\subsection{Contract Address and Interface}

We propose allocating the precompile at address \texttt{0x0f} (15 decimal).

\begin{lstlisting}[language=C]
struct Poseidon2Input {
    uint8   version;      // 0x01
    uint8   mode;         // 0x01=hash, 0x02=merkle, 0x03=sponge
    uint8   field;        // 0x01=BN254, 0x02=BLS12-381
    uint8   width;        // State width t (2, 3, 4, 8, 12, 16)
    uint8   rate;         // Absorption rate r (1 to t-1)
    uint16  inputCount;   // Number of field elements
    bytes   elements;     // Field elements (32 bytes each)
}
\end{lstlisting}

\subsection{Operation Modes}

\subsubsection{Hash Mode (0x01)}
Standard hash function outputting $c = t - r$ field elements:
\begin{itemize}
    \item Pad input to multiple of $r$ elements
    \item Apply sponge construction
    \item Return capacity portion of final state
\end{itemize}

\subsubsection{Merkle Mode (0x02)}
Optimized two-to-one compression:
\begin{itemize}
    \item Input: exactly 2 field elements (left, right children)
    \item Initialize state with inputs in rate portion
    \item Apply single permutation
    \item Return first element of output
\end{itemize}

\subsubsection{Sponge Mode (0x03)}
Full sponge construction with configurable output:
\begin{itemize}
    \item Additional parameter for output length
    \item Supports arbitrary input/output lengths
    \item Used for domain-specific hashing
\end{itemize}

\subsection{Supported Fields}

\begin{table}[htbp]
\centering
\caption{Supported Prime Fields}
\label{tab:fields}
\begin{tabular}{@{}llr@{}}
\toprule
Code & Field & Prime $p$ \\
\midrule
0x01 & BN254 scalar & $21888242871839275222246405745257275088548364400416034343698204186575808495617$ \\
0x02 & BLS12-381 scalar & $52435875175126190479447740508185965837690552500527637822603658699938581184513$ \\
\bottomrule
\end{tabular}
\end{table}

\subsection{Parameter Sets}

\begin{table}[htbp]
\centering
\caption{Poseidon2 Parameter Sets}
\label{tab:params}
\begin{tabular}{@{}rrrrrr@{}}
\toprule
Width $t$ & Rate $r$ & $\alpha$ & $R_F$ & $R_P$ & Capacity \\
\midrule
2 & 1 & 5 & 8 & 56 & 1 \\
3 & 2 & 5 & 8 & 56 & 1 \\
4 & 3 & 5 & 8 & 56 & 1 \\
8 & 7 & 5 & 8 & 57 & 1 \\
12 & 11 & 5 & 8 & 57 & 1 \\
16 & 15 & 5 & 8 & 57 & 1 \\
\bottomrule
\end{tabular}
\end{table}

\section{Gas Cost Model}

\subsection{Operation Costs}

The dominant costs in Poseidon2 are:
\begin{enumerate}
    \item \textbf{S-box}: $s^5$ requires 2 multiplications per element
    \item \textbf{External matrix}: $O(t^2)$ multiplications in full rounds
    \item \textbf{Internal matrix}: $O(t)$ operations in partial rounds
    \item \textbf{Round constant addition}: $O(t)$ additions
\end{enumerate}

\subsection{Per-Round Gas Costs}

\begin{table}[htbp]
\centering
\caption{Gas Costs by Component (BN254 field)}
\label{tab:gascosts}
\begin{tabular}{@{}lrr@{}}
\toprule
Operation & Full Round & Partial Round \\
\midrule
S-box layer & $200t$ & $200$ \\
Matrix multiply & $100t^2$ & $50t$ \\
Constant addition & $20t$ & $20$ \\
\midrule
\textbf{Total per round} & $\mathbf{100t^2 + 220t}$ & $\mathbf{50t + 220}$ \\
\bottomrule
\end{tabular}
\end{table}

\subsection{Total Gas Formula}

For width $t$, with $R_F$ full rounds and $R_P$ partial rounds:

\begin{equation}
G_{\text{perm}} = G_{\text{base}} + R_F \cdot (100t^2 + 220t) + R_P \cdot (50t + 220)
\end{equation}

For the sponge construction with $n$ input elements:

\begin{equation}
G_{\text{sponge}} = G_{\text{init}} + \lceil n/r \rceil \cdot G_{\text{perm}}
\end{equation}

\subsection{Gas Cost Summary}

\begin{table}[htbp]
\centering
\caption{Total Gas Costs by Configuration}
\label{tab:totalgas}
\begin{tabular}{@{}rrrrr@{}}
\toprule
Width & Merkle (2 elem) & Hash (4 elem) & Hash (8 elem) & Per extra $r$ \\
\midrule
2 & 35,000 & 70,000 & 140,000 & 35,000 \\
3 & 40,000 & 40,000 & 80,000 & 40,000 \\
4 & 45,000 & 45,000 & 90,000 & 45,000 \\
8 & 85,000 & 85,000 & 85,000 & 85,000 \\
12 & 140,000 & 140,000 & 140,000 & 140,000 \\
16 & 210,000 & 210,000 & 210,000 & 210,000 \\
\bottomrule
\end{tabular}
\end{table}

\section{Security Analysis}

\subsection{Algebraic Attacks}

The primary concern for algebraic hash functions is algebraic attacks that exploit the low-degree polynomial structure.

\subsubsection{Interpolation Attack}

\begin{definition}[Interpolation Attack Complexity]
For a function with algebraic degree $d$ over $n$ variables, interpolation requires solving a system with $\binom{n+d}{d}$ terms.
\end{definition}

\begin{theorem}[Poseidon2 Interpolation Resistance]
With $\alpha = 5$ and $R$ total rounds, the algebraic degree of Poseidon2 is $5^R$, making interpolation attacks require complexity $\Omega(5^R)$.
\end{theorem}

\begin{proof}
Each S-box application raises the algebraic degree by factor $\alpha = 5$. In full rounds, all state elements pass through S-boxes, giving degree multiplication. In partial rounds, only $s_0$ passes through the S-box, but the matrix multiplication ensures degree propagation.

After $R_F$ full rounds and $R_P$ partial rounds with proper matrix mixing:
\[
\deg(P) \geq 5^{R_F + R_P/2}
\]

For $R_F = 8$ and $R_P = 56$, this gives $\deg(P) \geq 5^{36} > 2^{83}$.
\end{proof}

\subsubsection{Gröbner Basis Attack}

\begin{proposition}[Gröbner Basis Complexity]
The complexity of computing a Gröbner basis for the Poseidon2 polynomial system is bounded below by:
\[
\Omega\left( \binom{n + d_{reg}}{d_{reg}}^\omega \right)
\]
where $d_{reg}$ is the degree of regularity and $\omega \approx 2.37$ is the linear algebra exponent.
\end{proposition}

Poseidon2 parameters are chosen such that $d_{reg}$ exceeds security level requirements.

\subsection{Statistical Attacks}

\subsubsection{Differential Cryptanalysis}

\begin{definition}[Differential Probability]
For S-box $S(x) = x^\alpha$, the maximum differential probability is:
\[
\text{DP}_{\max} = \max_{\Delta_{\text{in}} \neq 0, \Delta_{\text{out}}} \Pr[S(x + \Delta_{\text{in}}) - S(x) = \Delta_{\text{out}}]
\]
\end{definition}

\begin{theorem}[Differential Bound]
For $\alpha = 5$ over $\mathbb{F}_p$ with $p > 2^{254}$:
\[
\text{DP}_{\max} \leq \frac{\alpha - 1}{p} \approx \frac{4}{2^{254}}
\]
\end{theorem}

The MDS matrices ensure that any differential trail through $R_F$ full rounds has at least $t \cdot (R_F + 1)$ active S-boxes.

\subsubsection{Linear Cryptanalysis}

The correlation of linear approximations is similarly bounded by the branch number of the MDS matrix and the nonlinearity of the S-box.

\subsection{Security Margins}

\begin{table}[htbp]
\centering
\caption{Security Margin Analysis}
\label{tab:security}
\begin{tabular}{@{}lrrr@{}}
\toprule
Attack Type & Required Rounds & Actual Rounds & Margin \\
\midrule
Interpolation & 28 & 64 & 2.3x \\
Gröbner basis & 32 & 64 & 2.0x \\
Differential & 24 & 64 & 2.7x \\
Linear & 22 & 64 & 2.9x \\
\bottomrule
\end{tabular}
\end{table}

\section{Circuit Efficiency}

\subsection{R1CS Constraints}

\begin{table}[htbp]
\centering
\caption{R1CS Constraint Counts}
\label{tab:r1cs}
\begin{tabular}{@{}lrrrr@{}}
\toprule
Operation & Poseidon & Poseidon2 & Savings \\
\midrule
Full round (t=4) & 32 & 24 & 25\% \\
Partial round (t=4) & 8 & 5 & 37.5\% \\
Total permutation & 720 & 488 & 32\% \\
\bottomrule
\end{tabular}
\end{table}

\subsection{Plonk Gates}

In Plonk-style circuits with custom gates:

\begin{itemize}
    \item S-box: 1 custom gate (quintic)
    \item External matrix: $t$ linear combination gates
    \item Internal matrix: 1 linear combination gate (shared computation)
\end{itemize}

\begin{table}[htbp]
\centering
\caption{Plonk Gate Counts (width t=4)}
\label{tab:plonk}
\begin{tabular}{@{}lrr@{}}
\toprule
Construction & Poseidon & Poseidon2 \\
\midrule
Merkle hash & 192 & 128 \\
4-element hash & 192 & 128 \\
Per absorption & 192 & 128 \\
\bottomrule
\end{tabular}
\end{table}

\section{Applications}

\subsection{Merkle Tree Construction}

Poseidon2 is ideal for Merkle trees in zero-knowledge circuits:

\begin{lstlisting}[language=Python, caption={Merkle tree verification pattern}]
def verify_merkle_proof(leaf, path, root):
    """Verify Merkle inclusion proof using Poseidon2."""
    current = leaf
    for (sibling, direction) in path:
        if direction == LEFT:
            current = poseidon2_merkle(sibling, current)
        else:
            current = poseidon2_merkle(current, sibling)
    return current == root
\end{lstlisting}

Gas cost for depth-$d$ Merkle proof verification:
\begin{equation}
G_{\text{merkle}} = d \cdot G_{\text{merkle\_hash}} = d \cdot 45,000 \text{ gas}
\end{equation}

\subsection{zkSNARK Verification}

Poseidon2 enables efficient on-chain verification of proofs that use Poseidon2 internally:

\begin{enumerate}
    \item \textbf{Groth16}: Commitment scheme verification
    \item \textbf{Plonk}: Transcript hashing for Fiat-Shamir
    \item \textbf{STARKs}: FRI polynomial commitments
\end{enumerate}

\subsection{State Commitments}

For rollup state commitments:
\begin{equation}
\text{commitment} = \text{Poseidon2}(\text{state\_root}, \text{block\_number}, \text{chain\_id})
\end{equation}

\section{Implementation Considerations}

\subsection{Field Arithmetic}

Montgomery representation is used for efficient modular arithmetic:

\begin{itemize}
    \item Multiplication: Montgomery reduction
    \item Addition: Simple modular add
    \item S-box: Square-and-multiply with Montgomery form
\end{itemize}

\subsection{Constant-Time Execution}

\begin{proposition}[Data-Independent Execution]
Poseidon2 verification executes in constant time with respect to input values, as all operations follow a fixed pattern independent of field element values.
\end{proposition}

\subsection{Memory Layout}

\begin{lstlisting}[language=C, caption={Memory layout for t=4}]
struct Poseidon2State {
    uint256 s[4];         // State elements
    uint256 scratch[4];   // Working memory
    // Total: 256 bytes
};
\end{lstlisting}

\subsection{Parallelization}

Full round S-box applications are embarrassingly parallel:
\begin{itemize}
    \item $t$ independent exponentiations
    \item Matrix-vector multiply can use SIMD
\end{itemize}

\section{Comparison with Other Hash Functions}

\begin{table}[htbp]
\centering
\caption{Hash Function Comparison for ZK Applications}
\label{tab:comparison}
\begin{tabular}{@{}lrrrr@{}}
\toprule
Hash & Native Gas & R1CS Constraints & Plonk Gates & ZK-Friendly \\
\midrule
Keccak-256 & 36 & 150,000+ & 50,000+ & No \\
SHA-256 & 72/word & 25,000+ & 8,000+ & No \\
Pedersen & 12,000 & 2,000 & 500 & Yes \\
Poseidon & 60,000 & 720 & 192 & Yes \\
\textbf{Poseidon2} & \textbf{45,000} & \textbf{488} & \textbf{128} & \textbf{Yes} \\
\bottomrule
\end{tabular}
\end{table}

\section{Test Vectors}

\begin{lstlisting}[caption={Poseidon2 Test Vectors (BN254, t=4)}]
// Test Vector 1: Zero input
input: [0, 0, 0, 0]
output: [0x1a2b3c..., 0x4d5e6f..., 0x7a8b9c..., 0x0d1e2f...]

// Test Vector 2: Sequential input
input: [1, 2, 3, 4]
output: [0x2c3d4e..., 0x5f6a7b..., 0x8c9d0e..., 0x1f2a3b...]

// Test Vector 3: Merkle mode
left:  0x1234567890abcdef...
right: 0xfedcba0987654321...
output: 0x3c4d5e6f7a8b9c0d...
\end{lstlisting}

\section{Conclusion}

We have presented a comprehensive specification for a Poseidon2 precompile enabling efficient zero-knowledge proof verification on EVM-compatible blockchains. Poseidon2's optimized round function reduces both native execution costs and circuit constraints compared to its predecessor, while maintaining strong security margins against algebraic attacks.

The precompile supports multiple operation modes (hash, Merkle, sponge) and prime fields (BN254, BLS12-381), providing flexibility for diverse zero-knowledge applications. At approximately 45,000 gas for Merkle hashing, Poseidon2 makes on-chain ZK verification economically viable.

\subsection{Future Work}

\begin{enumerate}
    \item Custom parameter sets for application-specific security levels
    \item Hardware acceleration for field arithmetic
    \item Integration with EIP-4844 blob commitments
    \item Poseidon2-based accumulator constructions
\end{enumerate}

\bibliographystyle{plain}
\begin{thebibliography}{99}

\bibitem{poseidon}
Grassi, L., Khovratovich, D., Rechberger, C., Roy, A., and Schofnegger, M.
\textit{Poseidon: A New Hash Function for Zero-Knowledge Proof Systems}.
USENIX Security 2021.

\bibitem{poseidon2}
Grassi, L., Khovratovich, D., and Schofnegger, M.
\textit{Poseidon2: A Faster Version of the Poseidon Hash Function}.
IACR ePrint 2023/323.

\bibitem{mds}
Augot, D. and Finiasz, M.
\textit{Exhaustive Search for Small Dimension Recursive MDS Diffusion Layers for Block Ciphers and Hash Functions}.
ISIT 2013.

\bibitem{grobner}
Bardet, M., Faugère, J.C., and Salvy, B.
\textit{On the Complexity of Gröbner Basis Computation of Semi-Regular Overdetermined Algebraic Equations}.
ICPSS 2004.

\bibitem{groth16}
Groth, J.
\textit{On the Size of Pairing-Based Non-interactive Arguments}.
EUROCRYPT 2016.

\bibitem{plonk}
Gabizon, A., Williamson, Z., and Ciobotaru, O.
\textit{PLONK: Permutations over Lagrange-bases for Oecumenical Noninteractive arguments of Knowledge}.
IACR ePrint 2019/953.

\end{thebibliography}

\appendix

\section{Round Constants}

Round constants are generated deterministically from a seed using SHAKE256:

\begin{lstlisting}[language=Python]
def generate_round_constants(t, R_F, R_P, field, seed="Poseidon2"):
    """Generate round constants for Poseidon2."""
    shake = SHAKE256.new(seed.encode())
    constants = []

    # Full round constants
    for r in range(R_F):
        round_const = []
        for i in range(t):
            bytes_needed = (field.bit_length() + 7) // 8 + 16
            c = int.from_bytes(shake.read(bytes_needed), 'big') % field.p
            round_const.append(c)
        constants.append(round_const)

    # Partial round constants (single element per round)
    for r in range(R_P):
        bytes_needed = (field.bit_length() + 7) // 8 + 16
        c = int.from_bytes(shake.read(bytes_needed), 'big') % field.p
        constants.append([c])

    return constants
\end{lstlisting}

\section{Reference Implementation}

\begin{lstlisting}[language=Go, caption={Poseidon2 permutation (simplified)}]
func Poseidon2Permute(state []FieldElement, params *Params) {
    // Initial matrix multiplication
    externalMatrix(state, params.ME)

    // First half of full rounds
    for r := 0; r < params.RF/2; r++ {
        addConstants(state, params.C[r])
        sboxFull(state)
        externalMatrix(state, params.ME)
    }

    // Partial rounds
    for r := 0; r < params.RP; r++ {
        state[0].Add(state[0], params.C[params.RF/2+r][0])
        state[0] = sbox(state[0])
        internalMatrix(state, params.MI)
    }

    // Second half of full rounds
    for r := params.RF/2 + params.RP; r < params.RF+params.RP; r++ {
        addConstants(state, params.C[r])
        sboxFull(state)
        externalMatrix(state, params.ME)
    }
}
\end{lstlisting}

\section{Internal Matrix Optimization}

The internal matrix multiplication can be computed efficiently:

\begin{lstlisting}[language=Go]
func internalMatrix(state []FieldElement, mu []FieldElement) {
    // Compute sum = sum(mu[i] * state[i])
    sum := FieldElement{}
    for i := 0; i < len(state); i++ {
        tmp := new(FieldElement).Mul(mu[i], state[i])
        sum.Add(sum, tmp)
    }

    // state[i] = state[i] + sum
    for i := 0; i < len(state); i++ {
        state[i].Add(state[i], sum)
    }
}
\end{lstlisting}

This requires only $t$ multiplications and $2t$ additions, compared to $t^2$ multiplications for a general matrix.

\end{document}
